% ============================================================================
% Chapter 5: Protocol Implementation
% ============================================================================
\section{Protocol Implementation}
\label{sec:protocols}

Having analyzed the general architecture of Coroute, this chapter focuses directly on the concrete implementation of the supported protocols: HTTP/1.1, HTTP/2, WebSocket, and TLS. For each protocol, we will outline the key aspects of the implementation, including parsing, serialization, and the handling of specific edge cases.

The implementation of these protocols adheres to the principle of efficiency---minimizing data copying, utilizing zero-copy techniques wherever possible, and leveraging asynchronous processing via coroutines. Concurrently, the codebase is designed to remain readable and maintainable.

\subsection{HTTP/1.1 Parsing and Serialization}

HTTP/1.1 is a text-based protocol, which makes parsing relatively straightforward but requires careful attention to detail. The parser must handle various edge cases such as incomplete requests, invalid headers, and different encodings.

\subsubsection{Request Parsing}

HTTP/1.1 requests are parsed within the \texttt{parse\_request()} method. The process is asynchronous---data is read in chunks from the network until a complete request is accumulated. This is crucial for efficiency, as it allows the server to handle other requests while waiting for data from a slow client.

\begin{enumerate}
    \item Read headers until \texttt{\textbackslash r\textbackslash n\textbackslash r\textbackslash n} is encountered.
    \item Parse the request line (method, path, version).
    \item Parse individual headers.
    \item Read the body (if a \texttt{Content-Length} is specified).
\end{enumerate}

\begin{lstlisting}[language=C++, basicstyle=\small\ttfamily, caption={Parsing an HTTP request}]
Task<expected<Request, Error>> App::parse_request(net::Connection& conn) {
    constexpr size_t MAX_HEADER_SIZE = 8192;
    constexpr size_t MAX_BODY_SIZE = 10 * 1024 * 1024; // 10MB
    
    // Read headers in chunks
    auto buffer_ptr = buffer_pool_.acquire(MAX_HEADER_SIZE);
    auto& buffer = *buffer_ptr;
    
    size_t total_read = 0;
    size_t header_end_pos = std::string::npos;
    
    while (total_read < MAX_HEADER_SIZE) {
        auto result = co_await conn.async_read(
            buffer.data() + total_read, READ_CHUNK_SIZE);
        if (!result) co_return unexpected(result.error());
        
        total_read += *result;
        
        // Search for \r\n\r\n
        for (size_t i = search_start; i + 3 < total_read; ++i) {
            if (buffer[i] == '\r' && buffer[i+1] == '\n' && 
                buffer[i+2] == '\r' && buffer[i+3] == '\n') {
                header_end_pos = i + 4;
                break;
            }
        }
        if (header_end_pos != std::string::npos) break;
    }
    
    // Parse request line and headers...
    co_return req;
}
\end{lstlisting}

\subsubsection{URL Decoding}

URL addresses often contain special characters that are encoded using percent-encoding (for example, a space is encoded as \texttt{\%20}). Coroute automatically decodes URL parameters so that user code receives the clean values.

The \texttt{url\_decode} function handles three cases: percent-encoded characters (\texttt{\%XX}), the plus sign (which is interpreted as a space in query strings), and regular characters. The decoding is performed in-place for maximum efficiency.

\begin{lstlisting}[language=C++, basicstyle=\small\ttfamily, caption={URL decode function}]
static std::string url_decode(std::string_view str) {
    std::string result;
    result.reserve(str.size());
    
    for (size_t i = 0; i < str.size(); ++i) {
        if (str[i] == '%' && i + 2 < str.size()) {
            // Decode %XX
            char hex[3] = {str[i + 1], str[i + 2], '\0'};
            long val = std::strtol(hex, nullptr, 16);
            result += static_cast<char>(val);
            i += 2;
        } else if (str[i] == '+') {
            result += ' ';
        } else {
            result += str[i];
        }
    }
    return result;
}
\end{lstlisting}

\subsubsection{Keep-Alive Support}

HTTP/1.1 introduced the concept of persistent connections (keep-alive), which allows multiple HTTP requests to be sent over a single TCP connection. This represents a significant improvement over HTTP/1.0, where each request required a new TCP connection.

Coroute supports keep-alive with configurable parameters. By default, a connection can serve up to 100 requests before it is closed. There is also a timeout---if the client does not send a new request within 30 seconds, the connection is closed. These values can be configured according to the application's needs.

It is important to note that keep-alive is active by default in HTTP/1.1. The client must explicitly send \texttt{Connection: close} if it wants the connection to close after the request.

\begin{lstlisting}[language=C++, basicstyle=\small\ttfamily, caption={Keep-alive loop}]
constexpr size_t MAX_REQUESTS_PER_CONNECTION = 100;
constexpr auto KEEP_ALIVE_TIMEOUT = std::chrono::seconds(30);

size_t request_count = 0;
bool keep_alive = true;

while (conn->is_open() && keep_alive) {
    ++request_count;
    
    if (request_count > MAX_REQUESTS_PER_CONNECTION) break;
    
    auto req_result = co_await parse_request(*conn);
    // ... handle request ...
    
    keep_alive = req.keep_alive();
    
    // Set Connection header
    if (!keep_alive || request_count >= MAX_REQUESTS_PER_CONNECTION) {
        resp.set_header("Connection", "close");
    } else {
        resp.set_header("Connection", "keep-alive");
    }
}
\end{lstlisting}

\subsubsection{Zero-Copy File Transfer}

When serving static files, the traditional approach requires reading the file into user-space buffers and then writing to the socket. This imposes two distinct data copies---from the file system into the application buffer, and from the buffer out to the network stack.

Coroute utilizes zero-copy techniques, allowing the operating system to transfer data directly from the file system to the network stack without copying into user-space buffers. On Windows, this is achieved via \texttt{TransmitFile}, on Linux through \texttt{sendfile} or \texttt{splice}, and on macOS using \texttt{sendfile}.

Zero-copy proves exceptionally efficient for large files, where avoiding extra copies significantly improves performance while simultaneously reducing CPU load.

\begin{lstlisting}[language=C++, basicstyle=\small\ttfamily, caption={Zero-copy file serving}]
if (resp.has_file() && req.method() != HttpMethod::HEAD) {
    // Send headers first
    auto headers_data = resp.serialize_headers();
    co_await conn->async_write_all(headers_data.data(), headers_data.size());
    
    // Send file via zero-copy (TransmitFile/sendfile)
    const auto& file_info = resp.file_info();
    co_await send_file_zero_copy(
        *conn, file_info.path, file_info.offset, file_info.length);
}
\end{lstlisting}

\subsection{HTTP/2 Support}

HTTP/2 \cite{rfc7540} represents a significant evolution of the HTTP protocol. Unlike the text-based HTTP/1.1 counterpart, HTTP/2 is a binary protocol introducing stream multiplexing, header compression, and other advanced optimizations.

The HTTP/2 implementation in Coroute is comprehensive, encompassing full support for the protocol's core features. The server handles HTTP/2 connections conventionally over TLS (h2) alongside unencrypted architectures (h2c), although the latter remains rarely utilized in production.

\subsubsection{ALPN Negotiation}

Application-Layer Protocol Negotiation (ALPN) is a TLS extension authorizing clients and servers to negotiate their application protocol during the initial TLS handshake. This constitutes the standard mechanism separating HTTP/1.1 from HTTP/2 implementations during HTTPS connection attempts.

During TLS configuration, Coroute registers supported protocols in order of priority. Typically, HTTP/2 (``h2'') maintains higher priority over HTTP/1.1 due to its superior performance capabilities. Following the TLS handshake sequence, the server verifies the selected protocol and dispatches the connection toward the corresponding protocol handler natively.

\begin{lstlisting}[language=C++, basicstyle=\small\ttfamily, caption={ALPN negotiation}]
// Configure ALPN protocols
if (http2_enabled_) {
    tls_config.alpn_protocols = {"h2", "http/1.1"};
} else {
    tls_config.alpn_protocols = {"http/1.1"};
}

// After TLS handshake, check negotiated protocol
auto* tls_conn = dynamic_cast<net::TlsConnection*>(conn_result->get());
if (tls_conn) {
    auto proto = tls_conn->negotiated_protocol();
    if (proto && *proto == "h2") {
        // Create HTTP/2 connection
        auto h2_conn = std::make_shared<http2::Http2Connection>(
            std::move(*conn_result));
        handle_http2_connection(h2_conn).start_detached();
    }
}
\end{lstlisting}

\subsubsection{HTTP/2 Connection}

The HTTP/2 connection natively manages multiple concurrent streams internally:

\begin{lstlisting}[language=C++, basicstyle=\small\ttfamily, caption={HTTP/2 connection handler}]
h2_conn->set_handler([this](Request& r) -> Task<Response> {
    auto match = router_.match(r.method(), r.path());
    if (match) {
        r.set_route_params(std::move(match.params));
    }
    co_return co_await middleware_chain_.execute_or_not_found(
        r, match.handler);
});

Task<void> App::handle_http2_connection(
    std::shared_ptr<http2::Http2Connection> h2_conn) {
    active_connections_.fetch_add(1);
    
    try {
        co_await h2_conn->run();
    } catch (const std::exception& e) {
        std::cerr << "HTTP/2 error: " << e.what() << std::endl;
    }
    
    active_connections_.fetch_sub(1);
}
\end{lstlisting}

\subsubsection{h2c Upgrade}

HTTP/2 can also be utilized without TLS via an h2c upgrade:

\begin{lstlisting}[language=C++, basicstyle=\small\ttfamily, caption={h2c upgrade}]
Task<bool> App::try_http2_upgrade(
    std::unique_ptr<net::Connection>& conn, Request& req) {
    if (!http2_enabled_) co_return false;
    
    // Check for h2c upgrade request
    if (!http2::is_h2c_upgrade_request(req)) co_return false;
    
    // Upgrade the connection
    auto h2_conn = co_await http2::upgrade_to_http2(
        std::move(conn), req);
    if (!h2_conn) co_return true;
    
    (*h2_conn)->set_handler([this](Request& r) -> Task<Response> {
        // ... same handler as above ...
    });
    
    handle_http2_connection(*h2_conn).start_detached();
    co_return true;
}
\end{lstlisting}

\subsection{WebSocket Protocol}

WebSocket \cite{rfc6455} constitutes a protocol designed explicitly targeting real-time bidirectional communication. Unlike standard HTTP operations, where clients unilaterally initiate communication, WebSocket empowers the server to transmit data downstream toward clients asynchronously.

Coroute integrates comprehensive WebSocket support safely, including transparent upgrades originating from HTTP connections, parsing both text and binary message frameworks, ping/pong messaging validating keep-alive persistence natively, alongside graceful connection termination logic natively. The API is designed cleanly coupling intuitively alongside internal backend coroutine executions organically.

\subsubsection{Checking for WebSocket Upgrade}

A WebSocket connection begins as a standard HTTP request containing headers indicating a desire to upgrade. The server must verify whether the request contains all necessary headers before performing the upgrade.

\begin{lstlisting}[language=C++, basicstyle=\small\ttfamily, caption={WebSocket upgrade detection}]
bool is_websocket_upgrade(const Request& req) {
    auto upgrade = req.header("Upgrade");
    auto connection = req.header("Connection");
    auto key = req.header("Sec-WebSocket-Key");
    auto version = req.header("Sec-WebSocket-Version");
    
    return upgrade && connection && key && version &&
           *upgrade == "websocket" &&
           connection->find("Upgrade") != std::string::npos &&
           *version == "13";
}
\end{lstlisting}

\subsubsection{Calculating the Accept Key}

The WebSocket handshake requires computing an accept key:

\begin{lstlisting}[language=C++, basicstyle=\small\ttfamily, caption={WebSocket accept key}]
std::string compute_accept_key(std::string_view client_key) {
    // Magic GUID from RFC 6455
    constexpr auto MAGIC_GUID = "258EAFA5-E914-47DA-95CA-C5AB0DC85B11";
    
    std::string combined = std::string(client_key) + MAGIC_GUID;
    
    // SHA-1 hash
    auto hash = sha1(combined);
    
    // Base64 encode
    return base64_encode(hash);
}
\end{lstlisting}

\subsubsection{Upgrade Process}

\begin{lstlisting}[language=C++, basicstyle=\small\ttfamily, caption={WebSocket upgrade}]
Task<expected<std::unique_ptr<WebSocketConnection>, Error>> 
upgrade_to_websocket(std::unique_ptr<Connection> conn, const Request& req) {
    // Create 101 Switching Protocols response
    Response resp = create_upgrade_response(req);
    
    // Send upgrade response
    auto data = resp.serialize();
    auto result = co_await conn->async_write_all(data.data(), data.size());
    if (!result) {
        co_return unexpected(result.error());
    }
    
    // Wrap connection in WebSocketConnection
    co_return std::make_unique<WebSocketConnectionImpl>(std::move(conn));
}
\end{lstlisting}

\subsubsection{WebSocket Handler Registration}

\begin{lstlisting}[language=C++, basicstyle=\small\ttfamily, caption={WebSocket handler}]
app.websocket("/ws/chat", [](std::unique_ptr<WebSocketConnection> ws) 
    -> Task<void> {
    while (true) {
        auto msg = co_await ws->receive();
        if (!msg) break;  // Connection closed
        
        if (msg->type == WebSocketMessage::Text) {
            // Echo back
            co_await ws->send(msg->data);
        }
    }
});
\end{lstlisting}

\subsubsection{WebSocket Messages}

WebSocket communication utilizes frames:

\begin{lstlisting}[language=C++, basicstyle=\small\ttfamily, caption={WebSocket message types}]
struct WebSocketMessage {
    enum Type { Text, Binary, Close, Ping, Pong };
    
    Type type;
    std::string data;
};

class WebSocketConnection {
public:
    // Receive next message
    virtual Task<expected<WebSocketMessage, Error>> receive() = 0;
    
    // Send text message
    virtual Task<expected<void, Error>> send(std::string_view text) = 0;
    
    // Send binary message
    virtual Task<expected<void, Error>> send_binary(
        std::span<const std::byte> data) = 0;
    
    // Close connection
    virtual Task<void> close(uint16_t code = 1000) = 0;
};
\end{lstlisting}

\subsection{TLS Integration}

Transport Layer Security (TLS) forms a crucial necessity behind modern web application architectures. Coroute delivers full-stack robust TLS support seamlessly integrating with OpenSSL, recognized universally as industry-standard core cryptography.

The internal TLS setup spans intuitively configurable parameters guaranteeing highly flexible integration environments safely. Simpler environments only necessitate supplying valid paths towards standard certificates alongside private keys. More complex operational frameworks can implement client certificate verification loops, explicit cipher suite configurations, along with custom parameter sets.

Coroute supports standard TLS 1.2 implementations alongside prioritizing native TLS 1.3 capabilities whenever systemically available. ALPN negotiation organically executes natively, guaranteeing seamless HTTP/1.1 alongside HTTP/2 protocol selection.

\subsubsection{TLS Configuration}

Configuring TLS requires at a minimum a certificate and a private key. The certificate can be self-signed for development or issued by a Certificate Authority for production.

\begin{lstlisting}[language=C++, basicstyle=\small\ttfamily, caption={TLS configuration}]
struct TlsConfig {
    std::string cert_file;      // Server certificate
    std::string key_file;       // Private key
    std::string ca_file;        // CA certificate (optional)
    std::string chain_file;     // Certificate chain (optional)
    bool verify_client = false; // Client certificate verification
    
    // ALPN protocols
    std::vector<std::string> alpn_protocols;
};

// Enable TLS on the app
app.enable_tls({
    .cert_file = "server.crt",
    .key_file = "server.key",
    .alpn_protocols = {"h2", "http/1.1"}
});
\end{lstlisting}

\subsubsection{TLS Listener}

\begin{lstlisting}[language=C++, basicstyle=\small\ttfamily, caption={TLS listener}]
if (tls_enabled_ && tls_ctx_) {
    tls_listener_ = std::make_unique<net::TlsListener>(
        std::move(listener_), *tls_ctx_);
    
    [this]() -> Task<void> {
        while (!cancel_source_.is_cancelled()) {
            auto conn_result = co_await tls_listener_->accept();
            if (!conn_result) continue;
            
            // Check ALPN and dispatch to appropriate handler
            // ...
        }
    }().start_detached();
}
\end{lstlisting}
